
%%FILE blueprint/src/chapter/analysis_mathlib.tex
\chapter{Analysis in Mathlib}
\label{chap:analysis-mathlib}


% lemmas and definitions that are already in Mathlib, but might be useful to write as background is useful here

%%FILE blueprint/src/chapter/distributions.tex
%\documentclass[12pt,headings=small,paper=A4,DIV=calc]{article}

%\newcommand{\Norm}[1]{X_\omega}
%\newcommand{\Th}{\mathrm{Th}}
%\newcommand{\tH}{{\tilde H}}
%
%\newcommand{\pt}[1]{\overline{#1}}
%\usepackage[style=alphabetic]{biblatex}
%\newcommand{\ch}{}
%\addbibresource{TopoSeminar/form.bib}
%\newcommand{\detail}[1]{} % [#1]
%
%\include{preamble}
%\usepackage{mathabx,epsfig}
%\usepackage{mathabx}
%\def\acts{\mathrel{\reflectbox{$\righttoleftarrow$}}}
%\DeclareMathOperator{\Sph}{Sph}
%\DeclareMathOperator{\dt}{dt}
%\DeclareMathOperator{\KO}{KO}
%\DeclareMathOperator{\Spec}{Spec}
%\DeclareMathOperator{\Supp}{Supp}
%%\DeclareMathOperator{\Vect}{Vect}
%\renewcommand{\SO}{\mathrm{SO}}
%\renewcommand{\d}{\ \mathrm{d} }
%%\newtheorem*{ex}{Example}
%\newtheorem*{claim}{Claim}

\chapter{Distributions}
Laurent Schwartz introduced the notion of distributions two talk about generalized solutions to differential equations. For his work he got the fields medal.

\section{Space of Distributions}

\begin{definition}
	$\cD(\Omega) = C_c^\infty(\Omega)$ is the set of test functions together with a topology determined by its converging sequences : $\phi_n \to \phi$ in $\cD(\Omega)$ if
	\begin{itemize}
		\item there exists a compact subset $K \subset \Omega$ such that $\Supp(\phi_n) \subset K$.
		\item For all multiindices $\alpha$ we have $\partial_\alpha \phi_n \to \partial_\alpha \phi$ in uniformly.
	\end{itemize}
	Convention: if $x \in \R^d \setminus \Omega$, then $\phi(x) := 0$.\\
	$\cD'(\Omega)$ is the topological dual space, i.e. the space of continuous linear functionals $D(\Omega) \to \C$ with the weak-*-convergence, i.e. pointwise convergence.
\end{definition}
\begin{remark}
	A notion of converging sequence on a set $X$ is
	\begin{itemize}
		\item The constant function on $x$ converges to $x$
		\item A sequence converges to $x$ iff any subsequence has a subsequence converging to $x$.
	\end{itemize}
	Note, that any subsequence of a converging sequence converges to the same limit.
	It induces a closure operator on $X$ (for $\overline{A \cup B} \subset \bar A \cup \bar B$ you use, that if a sequence in $A \cup B$ converges to $x$, then it has a subsequence (converging to $x$) lying in $A$ or in $B$).
\end{remark}
% \begin{remark}
%     What do we really need from the space of test functions to build up our theory? It should be a subset of $C^\infty(\Omega)$ closed under convolution.
% \end{remark}
\begin{example}
	Every locally integrable function $f \in L^1_{loc}(\Omega)$ gives us a distribution $\Lambda f \in D'(\Omega)$ defined by
	\[
	\Lambda f (\phi) = \int_{\Omega} f(x) \phi(x) \dm x
	\]
	which is well-defined because $\phi$ has compact support.
\end{example}
\begin{example}
	Let $\mu$ be a Radon measure on $\Omega$ (or more generally a signed Borel measure which is finite on compact subsets of $\Omega$). Then it defines a distribution:
	\[
	T_\mu(\phi) = \int_{\Omega} \phi \dm \mu
	\]
\end{example}
\begin{proof}
	As Borel sets are $\mu$-measurable, every continuous function is $\mu$-measurable. Let $K := \Supp \phi \subset \Omega $. Because $\mu(K) < \infty$ and $\max \phi(K) < \infty$ it follows that $\phi$ is $\mu$-integrable. Linearity follows from linearity of the integral. If $\phi_j \to \phi_*$ unifomrly, then $\Supp \phi_j  \subset K$ for all $j$ and $T_\mu \phi_j \to \Lambda_\mu \phi_*$.
\end{proof}
We have the following important special case:
\begin{example}[Dirac-$\delta$]
	We have $\delta \in D'(\Omega)$ given by
	\[
	\delta(\phi) := \phi(0)
	\]
\end{example}

\subsection{Convolution}
\begin{notation}
	For $\phi \in D$ define $\phi^R \in D$ as $\phi^R(x) = \phi(-x)$ and for $x \in \R^d$ we have the shift $\tau_x(\phi) \in D$ given by $\tau_x(\phi)(y) = \phi(y - x)$ .
\end{notation}
\begin{example}{\label{ex:conv}}
	For $f \in L^1_{loc}(\Omega), g \in D(\Omega)$ we have
	\[
	(f * g)(x) = \Lambda f (\tau_x (\psi^R))
	\]
\end{example}
\begin{proposition}

	Let $F \in \cD'(\Omega) , \psi \in D$. The following two distributions coincide:
	\begin{enumerate}
		\item The distribution determined by the smooth function $x \mapsto F(\tau_x (\psi^R))$.
		\item The distribution $\phi \mapsto F(\psi^R * \phi)$.
	\end{enumerate}
	We write this distribution as $F * \psi$.
\end{proposition}

\begin{proof}
	$\zeta := \psi^R$
	The function $x \mapsto F(\tau_x(\psi^R))$ is smooth :
	\begin{itemize}
		\item It is continuous: If $x_n \to x$, then $\tau_x (\zeta) - \tau_{x_n}(\zeta) \to 0 $ uniformly and the same holds for partial derivatives. Now $F$ is continuous hence $F(\tau_{\bullet}(\psi^R))$ is continuous and $F$ preserves the difference quotients, hence it will be also smooth.
		\item The two distributions coincide:
		% Recall
		% \begin{align}
		%     (\zeta * \phi)(y) = (\Lambda \tau_y \psi) (\phi)
		% \end{align}
		%    Hence
		\begin{align*}
			F(\psi^R * \phi)  &= F \left(\int(\tau_x \psi^R(\bullet) \phi(x) \dm x) \right) \\
			&\overset{!}{=} \int F(\tau_x \psi^R(\bullet) \phi(x)) \dm x \\
			&= \int F(\tau_x \psi^R) \phi(x) \dm x \\
			&= \Lambda(F(\tau_{\bullet} \psi^R))(\phi)
		\end{align*}
		Where we are allowed to pull out the integral by \ref{lemma:distribCommInt}
	\end{itemize}
\end{proof}
\begin{lemma}{\label{lemma:distribCommInt}}
	Let $\phi \in C_c^{\infty}(\Omega \times \Omega)$. Then for any $F \in D'(\Omega)$ we have
	\[
	F \left(\int_\Omega \phi(x,\_) \dm x \right) = \int_\Omega F(\phi(x ,\_)) \dm x
	\]
\end{lemma}
\begin{proof}
	Consider $S_\varepsilon \in D$ defined by
	\[S_\varepsilon^\phi(y) = \varepsilon^d \sum_{n \in \Z^d} \phi(n \varepsilon,y)\]
	which is a finite sum as $\phi$ has compact support.
	Then for all $\phi \in C_c^\infty(\Omega \times \Omega)$ one has a limit in $\cD$
	\[
	\int_\Omega \phi(x,\_) \dm x = \lim_{\varepsilon \to 0} S^\phi_\varepsilon
	\]
	Hence by continuity of $F$
	\begin{align*}
		F(\int_\Omega \phi(x,\_) \dm x) &= F(\lim_{\varepsilon \to 0} S^\phi_\varepsilon)     \\
		&= \lim_{\varepsilon \to 0} F(S^{\phi}_\varepsilon)     \\
		&= \lim_{\varepsilon \to 0} \varepsilon^d \sum_{n \in \Z^d} F(\phi(n \varepsilon,\_))     \\
		&= \lim_{\varepsilon \to 0} S^{F \circ \phi}_\varepsilon \\
		&= \int_\Omega F(\phi(x,\_)) \dm x
	\end{align*}
\end{proof}
Using the first definition we learn the following two things
\begin{example}
	From the description \ref{ex:conv}: Writing $\Lambda f * g$ is unambiguous.
\end{example}
\begin{example}
	We have $\delta_0 * f = \Lambda f$ for all $f \in D$.
\end{example}
\begin{lemma}{\label{lemma:cont}}
	Convolution $F * \psi$ is continuous in both variables.
\end{lemma}
\begin{proof}
	Continuity in the distribution variable is clear by pointwise convergence. \\
	For the continuity in the test function variable one uses that convolution with a fixed test function is a continuous function $\cD \to \cD$ and distributions are continuous.
\end{proof}
\begin{proposition}{\label{proposition:deltaInClosure}}
	There exists a sequence $\psi_n \in C_c^{\infty}(\Omega)$ such that $\Lambda \psi_n \to \delta_0$ in $D'(\Omega)$.
\end{proposition}
\begin{proof}
	Fix some $\psi \in D$ with $\int \psi(x) \dm x = 1$. Define $\psi_n (x) := n^d \psi(nx)$. Then
	\[
	(\Lambda \psi_n - \delta_0) (\phi )  = \int n^d \psi (n x) \phi(x) \dm x - \phi(0) = \int \psi (x) \cdot (\phi(x / n) - \phi(0)) \dm x \to 0
	\]
	where we used that $\phi(x/n) - \phi(0) \to 0$ uniformly. \\
\end{proof}
The following is how we view $L^1_{loc}(\Omega) \subset D'(\Omega)$.
\begin{corollary}
	Let $f,g \in L^1_{loc}(\Omega)$ such that $\Lambda_f = \Lambda_g$. Then $f = g$ almost everywhere.
\end{corollary}
\begin{proof}
	We have $0 = \Lambda(f-g)(\tau_{\bullet} \psi_n^R) = \psi_n * (f - g) \to \delta_0 * (f-g) = f-g$ in $L^1_{loc}$, hence $f = g$ almost everywhere.
\end{proof}
\begin{example}
	$\delta$ is not a function! I.e. its not of the form $\Lambda f$ for some $f \in L^1_{loc}$
\end{example}
\begin{proof}
	We have $\Delta|_{\Omega \setminus \{0\}}= \Lambda 0$ so if it would be a function, then it would be zero almost everywhere, hence 0. But this is a contradiction.
\end{proof}
\begin{corollary}
	Then $C^\infty(\R^d)$ is dense in $\cD'(\R^d)$.
\end{corollary}
\begin{proof}
	We know by \ref{proposition:deltaInClosure} that there exists $\Lambda \psi_n \to \delta_0$ in $\cD'$.
	Let $F$ be a distribution on $\R^d$. Now setting $F_n := F * \psi_n^R$ , yields pointwise
	\[
	F_n(\phi) = F (\psi_n * \phi) \to F(\delta_0 * \phi) = F (\phi)
	\]
	by \ref{lemma:cont}
	hence $F_n \to F$ in $\cD'$.  %Then there exists a sequence of smooth functions such that $\Lambda f_n \to F$.

\end{proof}
\subsection{Derivatives}
For $\phi , \psi \in D$ we have
\[
\int_\Omega \partial^\alpha \phi \psi \dm x = (-1)^{|\alpha|} \int_\Omega \phi \partial^\alpha \psi \dm x
\]
This motivates the definition
\begin{definition}
	For a multiindex $\alpha$ and a distribution $F$ define the distribution
	\[\partial^\alpha F (\phi) = (-1)^{|\alpha|}(F \partial^\alpha \phi)\]
\end{definition}
\begin{remark}
	let $f \in L^1_{loc}(\Omega)$. If there exists some $f' \in L^1_{loc}(\Omega)$, such that $\Lambda f' = \partial^\alpha \Lambda  f$ as distributions, then we call $f'$ the {\bf weak derivative} of $f$ with respect to $\alpha$.
\end{remark}
\begin{proposition}
	For $F \in D', \phi \in D $ , We have
	\[\partial^\alpha (F * \phi) = (\partial^\alpha F) * \phi = F * \partial^\alpha \phi\]
\end{proposition}
\begin{proof}
	First note, that holds in the case where $F$ is a test function, so that we have ordinary convolution.
	Then check pointwise. After erasing the sign $(-1)^{|\alpha|}$
	\[
	F(\phi^R * \partial^\alpha \psi) = F(\partial^\alpha (\phi^R * \psi)) = F((\partial^\alpha \phi) * \psi)
	\]

\end{proof}
\subsection{Support}
The Support of a distribution $F$ is the complement of the largest open subset $U$, such that $F(\phi) = 0 \forall \phi \in D, \Supp(\phi) \subset U$.  This definition is unambiguous: If $F$ vanishes on each $U_i$ for some index set $I$ and $\phi \in D$ such that the compact set $\Supp \phi \subset U = \bigcup U_i$, we may assume that $I$ is finite and choose a partition of unity $\Supp \eta_i \subset U_i$ and $\sum \eta_i = 1$. Then $F(\phi) = \sum F(\phi \eta_i) = 0$.
\subsection{Tempered Distributions}
We no enlarge our test space to the schwartz space $\cS = \cS(\R^d)$ consisting of smooth functions that are rapidly decreasing at $\infty$ with all derivatives.
\begin{definition}
	Consider the increasing sequence of norms on $C^\infty(\Omega)$ defined by
	\[
	\|\phi\|_N = \sup \{ | x^\beta (\partial_x^\alpha \phi(x)) | \mid x \in \R^d , |\alpha|,\|\beta\|\le N\}
	\]
	\[
	\cS = \{\phi \in C^\infty(\R^d) \mid \|\phi\|_N < \infty \forall N\}
	\]
	with the obvious notion of convergence.
\end{definition}
\begin{lemma}
	We have a continuous inclusions $\cD \subset \cS$, hence $\cS' \subset \cD'(\R^d)$. \\
	Moreover, this inclusion is dense. Hence being tempered is a property of distributions. %and can be checked by asking that for any $\phi \in \cS$ and any sequence $\phi_n \in \cD \subset \cS$ converging to $\phi$ in $\cS$, the limit $\lim_{n \to \infty} F(\phi_n)$ exists.
\end{lemma}
\begin{lemma}{\label{lemma:slowlyInc}}
	Let $f \in L^1_{loc}(\R^d)$ such that there exists $N \ge 0$ with
	\[
	\int_{|x| < R} |f(x)| \dm x = O(R^N) , \text{ as } R \to \infty
	\]
	The distribution $\Lambda_f$ is tempered.
\end{lemma}

\begin{example}
	This condition holds for functions in $L^p(\R^d)$ for $p \in [1,\infty]$
\end{example}
\begin{lemma}
	If $F \in \cD'$ has compact support then it is tempered: Choose $\eta \in \cD$ such that $\eta|_U = 1$ on some neighborhood $U \supset \Supp(F)$. Then $F(\eta \phi) = F(\phi)$ for all $\phi \in \cD$ and $\phi \mapsto F (\eta \phi)$ defines a continuous functional on $\cS$. It does not depend on the choice of $\eta$, because given another such $\eta'$ and any $\phi \in \cS$ we have $\Supp((\eta - \eta') \cdot \phi) \cap \Supp(F) = \varnothing$.
\end{lemma}
Let $F$ denote a tempered distribution.
\begin{lemma}

	\begin{itemize}
		\item All $\partial^\alpha F$ are tempered.
		\item Let $\psi \in C^\infty$ be slowly increasing, i.e. for each $\alpha$ exists $N_\alpha$ such that $\partial_x^\alpha \psi(x) = O(|x|^{N_\alpha})$. Then $\psi F$, defined by $(\psi F)(\phi)(F(\psi \phi)$ is tempered.
	\end{itemize}
\end{lemma}
\begin{example}
	If $\psi \in \cS$, then $F(\tau_\bullet (\psi^R))$ is slowly increasing. The other formulation is still valid because $\cS$ is stable under $*$.
\end{example}
What is the point of the Schwartz space?
\begin{definition}
	The fourier transformation is a continuous bijection
	\begin{align*}
		\cS &\to \cS \\
		\phi &\mapsto \hat \phi  = (\xi \mapsto \int_{\R^d} \phi(x) e^{-2\pi i x \xi} \dm x)
	\end{align*}
\end{definition}

\begin{lemma}
	We have
	\[\Lambda_{\hat \psi}(\phi) = \int_{\R^d} \hat \psi (x) \phi(x) \dm x = \int_{\R^d} \psi(x) \hat \phi(x) \dm x = \Lambda_{\psi}(\hat \phi)\]
\end{lemma}
So its easy to define a compatible generalization to the tempered distributions:
\begin{definition}
	Define
	\[
	\hat F (\phi) = F (\hat \phi)
	\]
	and similarly for the inverse transform $f \mapsto \check{f}$.
\end{definition}
We automatically have the inversion theorem for distributions, because it holds for test functions.
\begin{example}
	If $1 \in S$ denotes the constant function at $1$, then
	\[
	\widehat {\delta} = \Lambda_1
	\]
	because
	\[
	\hat \delta (\phi) =\hat \phi(0) = \int 1 \phi(x) \dm x = \Lambda_1(\phi)
	\]
\end{example}

\section{Fundamental solutions}
In this chapter we may omit the $\Lambda$.
Fix a partial differential operator $L$
\[
L = \sum_{|\alpha|\le m } a_\alpha \partial^\alpha \text { on } \R^d
\]
with $a_\alpha \in \C$.
\begin{definition}
	A fundamental solution of $L$ is a distribution $F$ such that $L(F) = \delta$
\end{definition}
The reason why this is interesting:
\begin{lemma}
	The operator
	\[
	f \mapsto T(f) := F * f
	\]
	defines an inverse to $L$
\end{lemma}
\begin{proof}
	because
	\[
	\partial^\alpha (F * f) = (\partial^\alpha F) * f = F * (\partial^\alpha f)
	\]
	Summing over $\alpha$ gives
	\[
	L(F * f) = \delta * f  = F * L f
	\]
	Now recall $\delta * f = f$.
\end{proof}
\begin{definition}
	The characteristic polynomial of $L$ is
	\[P(\xi) = \sum_{|\alpha|\le m} a_\alpha (2 \pi i \xi)^\alpha\]
\end{definition}
It is defined it such a way, that $\widehat{L f} = P \cdot \hat f$.
So our hope is
\[F := \widecheck{1 / (P (\xi))} = \int \frac{1}{P(\xi)} e^{2 \pi i x \xi} \dm \xi\]
The problem is, that the zeros of $P$ result in difficulties to define $F$, even as a distribution.
\subsection{Laplacian}
If $L = \Delta = \sum_{i=1}^d \frac{\partial^2}{\partial^2 x_i}$
So $1 / P (\xi) = 1 / (-4 \pi^2 |\xi|^2) $ which lies in $L^1_{loc}$ for $d \ge 3$.
To calculate our distribution the following is helpful
\begin{theorem}
	For $\lambda > -d$, Let $H_\lambda$ be the tempered distribution associated to $|x|^\lambda \in L^1_{loc}$. \\
	If $-d  < \lambda < 0$  then
	\[
	\widehat {H_\lambda} = c_\lambda H_{-d-\lambda}
	\]
	with
	\[
	c_\lambda = \frac{\Gamma((d + \lambda)/2)}{\Gamma(\lambda / 2)} \pi^{-d/2-\lambda}
	\]
\end{theorem}
Set $\lambda:=-d+2$. Hence we can find an appropriate constant $C_d$ such that $F(x) = C_d |x|^{-d +2}$ is a fundamental solution, because $\hat {F}(\xi) = 1 / (-4\pi^2|\xi|^2)$. So writing out $\widehat{\Delta F} = 1$. \\

\begin{proposition}
	If $d = 2$, the function $F := 1/(2\pi) \log |x| \in L^1_{loc}$ is a fundamental solution of $\Delta$
\end{proposition}
\begin{proof} Sketch.
	One can actually compute $\hat F = -1 / (4\pi^2) \left [\frac{1}{|x|^2} \right] - c' \delta$ for some constant $c'$, where  $\left [\frac{1}{|x|^2} \right]$ is a distribution, that replaces the (non locally integrable) function $1  / |x|^2$ in an appropriate way:
	\[
	\left [\frac{1}{|x|^2} \right] = \int_{|x| \le 1} \frac{\phi(x) - \phi(0)}{|x|^2} \dm x + \int_{|x| > 1} \frac{\phi(x)}{|x|^2} \dm x
	\]
	Notice, that on the complement of zero, this distribution coincides with $1 / |x|^2$.
	\begin{notation}
		If $\phi \in C^\infty(\Omega)$ slowly increasing (i.e. all derivatives are bounded by polynomials), define $\phi F \in \cS'$ as $\phi F(\psi) = F(\phi \psi)$.
	\end{notation}
	\begin{align*}
		\widehat{ \Delta F} &= -4 \pi^2 |x|^2 \hat F\\
		&=  |x|^2 \left [\frac{1}{|x|^2} \right] - 4\pi^2 c' \underbrace{|x|^2 \delta}_{=0} \\
		&= 1
	\end{align*}
\end{proof}
\subsection{Heat operator}
$L = \frac{\partial}{\partial t} - \Delta_x$
taken over $(x,t) \in \R^{d+1} = \R^d \times \R$
i.e. we want to solve the homogeneous initial value problem
\[
\begin{cases}
	L (u) = 0 & , t > 0 \\
	u(x,0) = f(x) & , t = 0
\end{cases}
\]
for some initial value $f \in \cS$. \\
We have
\[
(\frac{\partial}{\partial t} \hat {\cH_t})(\xi) = \widehat{\frac{\partial}{\partial t} \cH_t}(\xi) = \widehat{\Delta_x \cH_t}(\xi) = -4\pi^2 |\xi|^2 \hat {\cH_t} (\xi)
\]
and this is obviously solved by $\cH_t = e^{-4\pi^2 |\xi|^2 t}$. We may call this the heat kernel
\[
\hat{\cH_t}(\xi) = e^{-4\pi^2 |\xi|^2 t}
\]
Note that for $t = 0$, we have $\hat{\cH_0} = 1$, hence $\cH_0  = \delta$, so $u(x,t) = (\cH * f)(x)$ solves the equation $L(u) = 0$ and $u(x,t) \to f(x)$ in $\cS$ as $t \to 0$.
\begin{remark}{\label{remark:approxId}}
	$\cH_t \to \delta$ in $\cS'$ as $t \to 0$ and $\int_{\R^d} \cH_t(x) \dm x = 1$ for all $t$
\end{remark}
Now define
\[
F(x,t) := \begin{cases}
	\cH_t(x) & , \ if  \ t > 0 , \\
	0 & , t \le 0
\end{cases}
\]
$F$ is locally integrable on $\R^{d+1}$ and it actually holds
\[
\int_{|t|\le R} \int_{\R^d} F(x,t) \dm x \dm t \le R
\]
so $F$ defines a tempered distribution by \ref{lemma:slowlyInc}.
\begin{theorem}
	$F$ is a fundamental solution of $L =  \frac{\partial}{\partial t} - \Delta_x$.
\end{theorem}
\begin{proof}
	Denote $L' = -  \frac{\partial}{\partial t} - \Delta_x$, then we have to see the last equation
	\[LF (\phi) = F(L'(\phi)) = \lim_{\varepsilon \to 0} \int_{t \ge \varepsilon} \int_{\R^d} F(x,t) ( -\frac{\partial}{\partial t} - \Delta_x) \phi(x,t) \dm x \dm t \overset{!}{=} \delta(\phi)
	\]
	Integration by parts
	\begin{align*}
		& \int_{t \ge \varepsilon} \int_{\R^d} F(x,t) ( -\frac{\partial}{\partial t} - \Delta_x) \phi(x,t) \dm x \dm t \\
		&= -  \int_{\R^d} \left(\int_{t \ge \varepsilon} \cH_t \frac{\partial}{\partial t} + (\Delta_x \cH_t) \phi \dm t \right) \dm x \\
		&= - \int_{\R^d} \left(\int_{t \ge \varepsilon} \cH_t \frac{\partial}{\partial t} + ( \frac{\partial}{\partial t} \cH_t) \phi \dm t \right ) \dm x \\
		&= \int_{\R^d} \cH_\varepsilon(x) \phi(x,\varepsilon) \dm x && \mid |\phi(x,\varepsilon) - \phi(x,0)| \le O(\varepsilon) \text{ uniformly in } x \\
		&= \int_{\R^d} \cH_\varepsilon(x) (\phi(x,0) + O(\varepsilon)) \dm x && \ref{remark:approxId}\\
		&\to \phi(0,0)
	\end{align*}
where in the last line we let $\varepsilon \to 0$
\end{proof}

\end{document}

%%FILE blueprint/src/chapter/interpolation.tex
\chapter{Interpolation}
\label{chap:interpolation}
% put text here

\section{Rietz-Thorin's Interpolation Theorem}
Rietz-Thorin's interpolation theorem is a powerful tool to study boundedness of linear operators between complex $L^p$ spaces.
Informally, it states that if a linear map $T$ is bounded as an operator $L^{p_0}\to L^{q_0}$ and as an operator $L^{p_1} \to L^{q_1}$, then
it must also be a bounded operator $L^p \to L^q$ whenever $\left(\frac{1}{p}, \frac{1}{q}\right)$ is a convex combination of $\left(\frac{1}{p_0}, \frac{1}{q_0}\right)$ and $\left(\frac{1}{p_1}, \frac{1}{q_1}\right)$.\\
Since simple functions are contained in all the $L^p$ spaces, and bounded linear operators are continuous, an equivalent formulation may be: given a bounded linear operator from simple functions to functions that are integrable on all sets of finite measure, if we know it can be extended to bounded linear
operators $L^{p_0} \to L^{q_0}$ and $L^{p_1} \to L^{q_1}$, then it can also be extended $L^p \to L^q$ with $p$ and $q$ as above.\\\\
Before we start, let us recall the maximum modulus principle from complex analysis.
There are various statements of this in Lean, see \href{https://leanprover-community.github.io/mathlib4_docs/Mathlib/Analysis/Complex/AbsMax.html}{the dedicated Mathlib page}.

\begin{theorem}
    \label{thm:maximum_modulus}
    \lean{Analysis.Complex.norm_eqOn_closure_of_isPreconnected_of_isMaxOn}
    \leanok
    Let $U$ be a connected open set in a complex normed space $E$. Let $f:E\to F$ be a function that
    complex differentiable on $U$ and continuous on $\bar{U}$.\\
    If $|f(z)|$ takes its maximum on a point $u\in U$, then it must be constant on $\bar{U}$.
\end{theorem}
\begin{proof}
    \leanok
    Already formalized in Mathlib, along with several variants.
\end{proof}


\begin{lemma}
    \label{lem:threelines}
    %\uses{}
    \lean{DiffContOnCl.norm_le_pow_mul_pow}
    \leanok
    Let $S$ be the strip $S:=\{z \in \C \ | \ 0 < \mathrm{Re} \, z < 1 \}$. Let $f: \overline{S} \to \C$
    be a function that is holomorphic on $S$ and continuous and bounded on $\overline{S}$.\\
    Assume $M_0, M_1$ are positive real numbers such that for all values of $y$ in $\R$, we have
    \[| \phi(iy) | \leq M_0 \ \qquad | \phi(1+iy) | \leq M_1 \]
    i.e., the absolute values of the function on the lines $\{\mathrm{Re} \, z = 0\}$ and $\{\mathrm{Re} \, z = 1\}$ are bounded by $M_0$ and $M_1$ respectively.\\
    Then, for all $0 \leq t \leq 1$ and for all real values of $y$, we have
    \[  | \phi(t + iy) | \leq M_0^{1-t} M_1^t\]
\end{lemma}
\begin{proof}
    % \uses{} % Put any results used in the proof but not in the statement here
    % Not sure what to put here exactly
    \leanok % uncomment if the lemma has been proven
    ~\\
    If $|\phi|$ is constant, everything holds trivially by setting $M_0$ and $M_1$ to be the value of $|\phi|$ at a point. Assume $|\phi|$ non-constant.\\
    \begin{itemize}
        \item{ \textbf{Case 1}: assume $M_0=M_1=1$, and $\sup_{0 \leq x \leq 1} | \phi(x+iy) | \to 0$ when $|y| \to \infty$.\\
        Let $M$ be the supremum of $|\phi(z)|$ on $\bar{S}$. Since the function is non-constant, we have $M>0$.\\
        Let $\{z_n\}$ be a sequence of points in $S$ such that $|\phi(z_n)|$ converges to $M$.\\
        Since we assumed the absolute value of $\phi$ goes to zero as $|y|$ goes to infinity, all points where $|\phi(z)|>M-\epsilon$ must be in some rectangle around zero, i.e. the sequence ${z_n}$ must be bounded.\\
        Hence, there must be a converging subsequence of ${z_n}$ to a point $z^* \in \bar{S}$.\\
        By maximum modulus principle, $z^*$ must be on the boundary $\delta S$, so it must have real part $0$ or $1$.
        Hence, by assumption, $|\phi(z^*)|\leq 1$, and by construction $|\phi(z)|\leq 1$ for all $z \in \bar{S}$, which is what we wanted to show.
        }
        \item{ \textbf{Case 2}: only assume $M_0 = M_1 = 1$.\\
        For $\epsilon>0$, define
        \[ \phi_{\epsilon} (z) = \phi(z) e^{\epsilon (z^2 - 1)} \]
        If the real part of $z$ is $0$, then $z=iy$ and
        \[ | \phi_{\epsilon} (z) | = |\phi(z) e^{\epsilon (-y^2-1)}| \leq |\phi(z)| \cdot 1 \leq 1\]
        If the real part of $z$ is $1$, then $z=1+iy$ and
        \[ | \phi_{\epsilon} (z) | = |\phi(z) e^{\epsilon (1 -y^2 + 2iy -1)}| = |\phi(z) e^{\epsilon (-y^2 + 2iy)}| = |\phi(z) e^{\epsilon (-y^2)}| \leq |\phi(z)| \cdot 1 \leq 1 \]
        Moreover,
        \[ | \phi_{\epsilon} (x+iy)| \leq |\phi(x+iy)| \cdot |e^{\epsilon(z^2-1)}| = |\phi(x+iy)| \cdot |e^{\epsilon(x^2-1-y^2+2ixy)}| = |\phi(x+iy)| \cdot |e^{\epsilon(x^2-1-y^2)}| \]
        Hence, for $0\leq x \leq 1$ and $|y| \to \infty$, we have that both factors go to zero.\\
        Thus $\phi_{\epsilon}$ satisfies the hypotheses of case 1, so $|\phi_{\epsilon}| \leq 1$ on the whole strip.\\
        Now, we have pointwise that
        \[\lim_{\epsilon \to 0} \phi_{\epsilon}(z) = \lim_{\epsilon \to 0} \phi(z) e^{\epsilon(z^2-1)} = \phi(z)\]
        Hence, for $\epsilon \to 0$, we have $|\phi_{\epsilon}(z)| \to |\phi(z)|$.
        Thus,
        \[ |\phi(z)| = \lim_{\epsilon \to 0} |\phi_{\epsilon}(z)| \leq 1 \]
        which is what we wanted to show.
        }
        \item{ \textbf{General case}\\
        If $M_0$ and $M_1$ are any two positive real numbers, define
        \[ \tilde{\phi} (z) = M_0 ^{z-1} M_1^{-z} \phi(z) \]
        Recall that, for $a\in \R\setminus\{0\}$, we have
        \[ |a^{b+ic}| = |a^b| \]
        Hence, if the real part of $z$ is $0$, we have
        \[ |\tilde{\phi} (z)| \leq |M_0^{-1}| \cdot |M_1^0| \cdot |\phi(z)| \leq \frac{1}{M_0} \cdot M_0 = 1 \]
        And if the real part of $z$ is $1$, we have
        \[ |\tilde{\phi} (z)| \leq |M_0^0| \cdot |M_1^{-1}| \cdot |\phi(z)| \leq \frac{1}{M_1} \cdot M_1 = 1 \]
        From the previous case, we obtain that for arbitrary $z$ in the strip
        \[ |\tilde{\phi} (z)| \leq 1 \]
        Now, write $z= t + iy$ and unroll the definition of $\tilde{\phi}$ to obtain
        \[ |M_0^{t-1+y} M_1^{-t-iy} \phi(z)| \leq 1 \]
        The left-hand side is equal to
        \[ M_0^{t-1} M_1^{-t} |\phi(z)| \]
        So we obtain
        \[|\phi(z)| \leq M_0^{1-t} M_1^t\]
        which is exactly what we wanted.
        }
    \end{itemize}
\end{proof}


\begin{lemma}
  \label{lem:snorm_eq_sSup_snorm}
  %\uses{}
  \lean{MeasureTheory.snorm_eq_sSup_snorm}
  \leanok
  Let $p$ and $q$ be real conjugate exponents. Let $f$ be measurable. Then
  \[ \|f\|_{L^q} = \sup_{\|g\|_{L^p} \leq 1, \ g \text{ simple}} \| fg \|_{L^1}.\]

  In particular, if the right hand side formula is finite, $f\in L^q$.
  \end{lemma}
\begin{proof}
  % \uses{} % Put any results used in the proof but not in the statement here
  \leanok % uncomment if the lemma has been proven
  That $$\sup_{\|g\|_{L^p} \leq 1, \ g \text{ simple}} \| fg \|_{L^1} \le \|f\|_{L^q}$$
  follows from Hölder's inequality.

  The other direction is trivial when $\|f\|_{L^q}=0$. Suppose $\|f\|_{L^q}\ne 0$.
  Set $$g = \cfrac{|f|^{p-1}}{\|f\|_{L^p}^{p-1}}.$$
  Then $\|g\|_{L^p} = 1$.
  \begin{align*}
    \|f\|_{L^q} &= \| fg \|_{L^1} \\
    &= \sup_{n \in \mathbb N} \| fg_n \|_{L^1} \tag{1}\\
    &\le \sup_{\|g\|_{L^p} \leq 1, \ g \text{ simple}} \| fg \|_{L^1} \tag{2}
  \end{align*}

  In (1) and (2), $g_n$ is a monotone sequence of simple function approximating $g$ from below,
  whose existence is basic real analysis.
\end{proof}

\begin{lemma}
  \label{lem:snormEssSup_eq_sSup_snorm}
  %\uses{}
  \lean{MeasureTheory.snormEssSup_eq_sSup_snorm}
  \leanok
  Let $f$ be measurable and the measure $\mu$ be $\sigma$-finite. Then
  \[ \|f\|_{L^\infty} = \sup_{\|g\|_{L^1} \leq 1, \ g \text{ simple}} \| fg \|_{L^1}.\]
\end{lemma}

\begin{proof}
  \leanok
  That $$\sup_{\|g\|_{L^1} \leq 1, \ g \text{ simple}} \| fg \|_{L^1} \le \|f\|_{L^\infty}$$
  follows from Hölder's inequality.

  Suppose $M:= \sup_{\|g\|_{L^1} \leq 1, \ g \text{ simple}} \| fg \|_{L^1} < \|f\|_{L^\infty}$.
  Then $\{x | |f(x)| \ge M \}$ has positive measure. Since $\mu$ is $\sigma$-finite, we have a subset $B \subset \{x | |f(x)| \ge M \}$ with finite positive measure
  by \href{https://leanprover-community.github.io/mathlib4_docs/Mathlib/MeasureTheory/Measure/Typeclasses.html#MeasureTheory.Measure.exists_subset_measure_lt_top}{a classical lemma.} Now set
  $$h := \mu(B)^{-1} \chi_B.$$
  Then we have
  \begin{align*}
    M &= \int_B \mu (B)^{-1} M d\mu \\
    &< \int_B \mu (B) |f| d\mu \\
    &= \|fh\|_{L^1}.
  \end{align*}
  But $\|h\|=1$ and $h$ is simple, so $\|fh\|_{L^1} \le M$, contradiction.

\end{proof}





\begin{comment}
      (Note: This is used by Stein-Shakarchi to prove the dual of $L^p$ is $L^q$, so it may be already formalized by the time
      we get to this point?)\\\\
      Approximate $g$ by simple functions from below, i.e. consider a sequence $\{g_n\}_{n\in \bbn}$ of simple functions
      such that for each $x$ we have $|g_n (x)| \leq |g(x)|$ and the $g_n$ converge to $g$ pointwise.\\
      If $p>1$ (thus $q<\infty$), consider
      \[ f_n(x) = |g_n(x)|^{q-1} \mathrm{sign}(g(x)) \cdot \frac{1}{||g_n||^{q-1}_{L^q}} \]
      Observe that $||f_n||_{L^p}=1$. For this, first note that $p(q-1)=q$ since
      \[\frac{1}{p} = 1- \frac{1}{q}=\frac{q-1}{q} \ \Rightarrow \ q = p(q-1)  \]
      and then write
      \[||f_n||^p_{L^p} = \int \left|\frac{1}{||g_n||_{L^q}^{q-1}} \cdot |g_n(x)|^{q-1} \mathrm{sign} g(x) \right|^p = \frac{1}{||g_n||_{L^q}^{p(q-1)}} \int  |g_n(x)|^{p(q-1)} = \frac{1}{||g_n||_{L^q}^q} ||g_n||_{L^q}^q = 1  \]
      So by assumption
      \[ \left|\int f_n g\right| \leq M \]
      And by direct computation
      \[ \int f_n g = \int \]

      HOLE HERE
      Then, we conclude by Fatou's Lemma

      If $p=1$, since the measure is $\sigma$-finite, write $X$ as an increasing union of subsets $\{E_n\}$ and take
      \[ f_n (x) = \frac{1}{\mu(E_n)} \mathrm{sign}(g(x)) \chi_{E_n} (x) \]
      Then $||f_n(x)||_{L^1} = \frac{1}{\mu(E_n)} \cdot \mu(E_n) = 1$.

      Now, it is easy to show for both cases that $||g||\geq M$ by H\"older's inequality: since $M$ is the supremum of the absolute value of $\int fg$ for our possible $f$,
      fix an $\epsilon>0$ and choose an $f$ such that
      \[\left| \int fg \right| \geq M-\epsilon \]
      so that
      \[ M- \epsilon \leq ||fg||_{L^1} \leq ||f||_{L^p} ||g||_{L^q} \leq ||g||_{L^q} \]
      Since $\epsilon$ is arbitrarily small, this gives
      \[ ||g||_{L^q} \geq M \]
\end{comment}

As a last step towards proving the theorem, let us recall a consequence of Hölder's inequality, which will only really be substantial in a corner case of our proof.
\begin{lemma}
  \label{lem:hoelder'}
  %\uses{}
  %\lean{}
  \leanok
  Let $(X, \mu)$ be a measure space and $0< p_0 <p_1\leq \infty$, $0 \le p \le \infty$. Assume we have $t \ge 0$ such that
  \[ \frac{1}{p} = \frac{1-t}{p_0} + \frac{t}{p_1} \]
  Let $f : X \to \mathbb C$ be a measurable function. Then
  \[ \| f \|_{L^p} \leq \| f \|^{1-t}_{L^{p_0}} \|f\|^t_{L^{p_1}} \]
  In particular, the $f$ in $L^{p_0}(X) \cap L^{p_1}(X)$, then $f$ is in $L^p(X)$.

  \end{lemma}
    \begin{proof}
    %\uses{} % Put any results used in the proof but not in the statement here
    \leanok % uncomment if the lemma has been proven
    % ``Exercise''/ To be filled later/ Possibly already in Mathlib

    This is just a version of Hölder's inequality, but in order to apply it, we should rule out some trivial cases.

    First note that the assumption guarantees that $0 < p_0 \le p \le p_1 \le \infty$ and $0\le t \le 1$.

    When $t=0$ or $t=1$, the inequality holds trivially. Hence we may assume $t \ne 0$ and $t \ne 1$. In this case $p < p_1 \le \infty$.

    Now we can rearrange the equality into
    $$ \cfrac{1}{\cfrac{p_0}{(1-t)p}} + \cfrac{1}{\cfrac{p_1}{tp}} = 1.$$

    Using Hölder's inequality, we have
    \begin{align*}
      \|f\|_{L^p} &= (\int |f|^p d\mu)^{1/p} \\
      & = (\int |f|^{(1-t)p} |f|^{tp} d\mu)^{1/p} \\
      & \le ((\int |f|^{(1-t)p \frac{p_0}{(1-t)p}} d\mu)^{\frac{(1-t)p}{p_0}} (\int |f|^{tp \frac{p_1}{tp}} d\mu)^{\frac{tp}{p_1}})^{1/p} \\
      & = (\int |f|^{p_0} d\mu)^{\frac{(1-t)}{p_0}} (\int |f|^{p_1} d\mu)^{\frac{tp}{p_1}} \\
      & = \| f \|^{1-t}_{L^{p_0}} \|f\|^t_{L^{p_1}}
    \end{align*}

  \end{proof}




\begin{theorem}
  \label{thm:riesz_interpolation}
  % \uses{}
  % \lean{}
  %\leanok
  Let $(X, \mu)$ and $(Y, \nu)$ be measure spaces and consider all $L^p$ spaces to be complex valued.\\
  Suppose $T$ is a linear map $L^{p_0}(X) + L^{p_1}(X) \to L^{q_0}(Y) + L^{q_1}(Y)$ that restricts to
  bounded operators $L^{p_0} \to L^{q_0}$ and $L^{p_1} \to L^{q_1}$. Let $M_0, M_1$ be the respective bounds, i.e.,
  \[ \begin{cases} ||Tf||_{L^{q_0}} \leq M_0 ||f||_{L^{p_0}} \\ ||Tf||_{L^{q_1}} \leq M_1 ||f||_{L^{p_1}}  \end{cases}\]
  Then, for any pair $(p, q)$ such that there is a $t$ in $[0,1]$ for which
  \[\frac{1}{p} = \frac{1-t}{p_0} + \frac{t}{p_1} \qquad \frac{1}{q} = \frac{1-t}{q_0} + \frac{t}{q_1} \]
  we have that the operator is bounded $L^p \to L^q$, and in particular
  \[ ||Tf||_{L^q} \leq M_0^{1-t} M_1^t ||f||_{L^p} \]
\end{theorem}
\begin{proof}
  \uses{lem:lp_of_sup_norm, lem:threelines} % Put any results used in the proof but not in the statement here
  % \leanok % uncomment if the lemma has been proven
  For a valid choice of $p,q$, note that we both need to show $Tf$ is in $L^q$ and a bound on the
  $L^q$ norm of $Tf$.
  Let $q'$ be the conjugate exponent of $q$. By Lemma \ref{lem:lp_of_sup_norm} for $Tf$, we need a bound of the form
  \[ \sup_{||g||_{L^{q'}} \leq 1, \ \ g \text{ simple}} \left| \int T(f) g \right| \leq M ||f||_{L^p} \]
  \begin{itemize}
    \item{\textbf{Case 1}: assume $p<\infty$ and $q>1$, and let $q', q_0^1, q_1'$ be the conjugate exponents of $q, q_0, q_1$ respectively.
    \begin{itemize}
        \item{\textbf{Subcase a:} assume $f=\sum_i a_i \chi_{E_i}$ is a simple function (finite sum with $E_i$ disjoint of finite measure).\\
        Let $g=\sum_j b_j \chi_{F_j}$ be a simple function.
        By writing $f$ as $||f||_{L^p} \cdot \frac{f}{||f||_{L^p}}$ and using linearity of $T$ and integrals, it suffices to prove the above inequality when $||f||_{L^p}=1$.\\
        We want to apply the three lines lemma to an appropriate function. Define
        \[ \gamma(z) = p \left(\frac{1-z}{p_0} + \frac{z}{p_1} \right) \qquad f_z = |f|^{\gamma(z)} \cdot \frac{f}{|f|} \]
        \[ \delta(z) = q' \left(\frac{1-z}{q'_0} + \frac{z}{q'_1} \right) \qquad g_z = |g|^{\delta(z)} \cdot \frac{g}{|g|} \]
        Observe that for $t$ as in the statement of the theorem, we have by definition that $\gamma(t)=1$, hence $f_t=f$.\\
        Moreover, if $\mathrm{Re}(z)=0$, we have that $\mathrm{Re} (\gamma(z)) = \frac{p}{p_0}$, and hence
        \[ ||f_z||_{L^{p_0}} = \left(\int |f_z|^{p_0}\right)^{\frac{1}{p_0}} = \left(\int | |f|^{\gamma(z)} |^{p_0}\right)^{\frac{1}{p_0}} = \left(\int |f| ^{\frac{p}{p_0} \cdot p_0}\right)^{\frac{1}{p_0}} = \left(||f||_{L^p}^p\right)^{\frac{1}{p_0}} = 1^{\frac{1}{p_0}} = 1  \]
        If $\mathrm{Re}(z)=1$, we have $\mathrm{Re} (\gamma(z)) = \frac{p}{p_1}$, and the exact same computation replacing $p_0$ with $p_1$ now shows that $||f_z||_{L^{p_1}}=1$.\\
        Similarly, one shows that
        \[ g_t=g \qquad ||g_z||_{L^{q'_0}} = 1 \text{ if } \mathrm{Re}(z)=0 \qquad ||g_z||_{L^{q'_1}} = 1 \text{ if } \mathrm{Re}(z)=1 \]
        Now, we want to apply the three lines lemma to the function
        \[\phi(z) := \int (Tf_z) g_z \]
        Since $f$ and $g$ are simple and given by the expressions above, we can explicitly write $f_z$ and $g_z$ as
        \[f_z=\sum_i |a_i|^{\gamma(z)} \frac{a_i}{|a_i|} \chi_{E_i} \qquad g_z=\sum_j |b_j|^{\delta(z)} \frac{b_j}{|b_j|} \chi_{F_j} \]
        (here, we use that the $E_i$ (respectively $F_j$) are disjoint, so for every point in the domain there is at most one of the $E_i$ covering it).\\
        So, expanding everything by linearity of $T$ and integrals, we obtain
        \[ \phi(z) = \sum_{i,j} |a_i|^{\gamma(z)} \frac{a_i}{|a_i|} |b_j|^{\delta(z)} \frac{b_j}{|b_j|} \int T(\chi_{E_i}) \chi_{F_j} \]
        This only depends on $z$ holomorphically in terms of the exponents of the $|a_i|$ and $|b_j|$, so it is a holomorphic function on the strip $S$ in the three lines lemma and it is continuous on $S$.\\
        It it also bounded on $\bar{S}$. In fact, we wrote $\phi$ as a finite sum, so we only need to show each summand is bounded.
        Since the real part of $z$ is between $0$ and $1$, the terms $|a_i|^{\gamma(z)}$ and $|b_j|^{\delta(z)}$ have bounded norms.
        Finally, recall that H\"older's inequality states that $||fg||_1 \leq ||f||_p ||g||_1$ for conjugate exponents $p,q$.
        Hence,
        \[ \left| \int T(\chi_{E_i}) \chi_{F_j} \right| \leq ||T(\chi_{E_i}) \chi_{F_j}||_{L^1} \leq ||T(\chi_{E_i})||_{L^{p_0}} \cdot ||\chi_{F_j}||_{p'_0} \leq M_0 \mu(E_i) \mu(F_j) \]
        which is bounded. Moreover, if $\mathrm{Re}(z)=0$, since $||f_z||_{L^{p_0}}=||g_z||_{L^{q'_0}}=1$, we have
        \[ | \phi(z) | \leq \int |(Tf_z) g_z| \leq ||Tf_z||_{L^{q_0}} \cdot ||g_z||_{L^{q'_0}} \leq M_0 \cdot ||g_z||_{L^{q'_0}} \leq M_0 \]
        Similarly, if $\mathrm{Re}(z)=1$, we obtain
        \[ | \phi(z) | \leq \int |(Tf_z) g_z| \leq ||Tf_z||_{L^{q_1}} \cdot ||g_z||_{L^{q'_1}} \leq M_1 \cdot ||g_z||_{L^{q'_1}} \leq M_1 \]
        Thus, applying the three lines lemma to $\phi(z)$ yields that
        \[ |\phi(t+yi)| \leq M_0^{1-t} M_1^t \]
        In particular, this holds for $y=0$, but now
        \[ \phi(t) = \int (Tf_t) g_t = \int (Tf) g \]
        So we have
        \[ \left| \int (Tf) g \right| \leq M_0^{1-t} M_1^t \]
        which is exactly what we wanted to show.
        }
        \item{\textbf{Subcase b:} Now, let $f$ be any function in $L^p$. By density of simple functions, approximate $f$ by a sequence ${f_n}$ of simple functions with $||f_n - f||_{L^p} \to 0$.\\
        By the previous case, we have $||Tf_n||_{L^q} \leq M ||f_n||_{L^p}$. In particular, the sequence $\{Tf_n\}$ is Cauchy in $L^q$, since
        \[ ||Tf_m - Tf_n||_{L^q} = ||T(f_m-f_n)||_{L^q} \leq M ||f_m - f_n||_{L^p} \]
        and the original sequence is Cauchy. By completeness, the $\{Tf_n\}$ converge in $L^q$, in particular the $L^q$ norm of the limit is the limit of the $L^q$ norms, which is less than $M ||f||_{L^p}$.
        Hence, it suffices to show that the sequence $\{Tf_n\}$ converges almost everywhere to $Tf$.\\
        Write $f= f^U + f^L$ with
        \[ f^U := \begin{cases} f(x) & \text{if } |f(x)|\geq 1 \\ 0 & \text{otherwise} \end{cases} \qquad f^L := \begin{cases} f(x) & \text{if } |f(x)|< 1 \\ 0 & \text{otherwise} \end{cases}  \]
        and similarly $f_n = f_n^U + f_n^L$.\\
        Modulo reordering them, assume $p_0 \leq p_1$, so we have $p_0 \leq p \leq p_1$. Since $f\in L^p$, $f^U$ must be in $L^{p_0}$ and $f^L$ in $L^{p_1}$.
        Similarly, since $f_n \to f$ in $L^p$, we have $f_n^U \to f^U$ in $L^{p_0}$ and $f_n^L \to f^L$ in $L^{p_1}$.\\
        By the assumptions of boundedness of $L$
        \[ Tf_n^U \to Tf^U \ \text{ in } L^{q_0} \qquad Tf_n^L \to Tf^L \ \text{ in } L^{q_1} \]
        Modulo extracting subsequences, we can assume that the convergence is almost everywhere, so that almost everywhere
        \[ Tf_n (x) = Tf_n^U (x) + Tf_n^L(x) \to Tf^U (x) + Tf(x) = Tf (x)\]
        which is what we wanted to show.
        }
    \end{itemize}
    }
    \item{\textbf{Case 2:} $p=\infty$ or $q=1$.\\
    If $p=\infty$, we must also have $p_0=p_1=\infty$, thus we have
    \[ ||Tf||_{L^{q_0}} \leq M_0 ||f||_{L^{\infty}} \qquad ||Tf||_{L^{q_1}} \leq M_1 ||f||_{L^{\infty}} \]
    Applying Lemma \ref{lem:hoelder'} with $Tf, q, q_0, q_1$, we obtain
    \[||Tf||_{L^q} \leq ||Tf||_{L^{q_0}}^{1-t} ||Tf||_{L^{q_1}}^{t} \leq  M_0^{1-t} M_1^t ||f||_{L^{\infty}} \]
    which is what we wanted.\\\\
    If $p<\infty$ and $q=1$, then (since they must be at least $1$ by definition of $L^p$ spaces) we have that $q_0$ and $q_1$ must also both be $1$ (for example, since $\frac{1}{q}$ is a convex combination of the other two reciprocals, the largest one must be $1$, and from that rearranging terms shows the other one is $1$).
    In this case, take $g_z=g$ for all $z$ and repeat the proof above (note:isn't this what already happens if we do not consider this case separately?).
    }
  \end{itemize}
\end{proof}

  \section{Applications of Rietz-Thorin's Interpolation Theorem}

  \subsection{Hausdorff-Young inequalities}

  \begin{lemma}
    \label{lem:hausdorff_young}
    % \uses{}
    % \lean{}
    %\leanok
    Let $X=[0,2\pi]$ with normalized Lebesgue measure $\frac{d\theta}{2\pi}$ and let $Y=\bbz$ with counting measure.\\
    Consider the operator $T: f \mapsto \{a_n\}_{n\in \bbz}$ where
    \[ a_n = \frac{1}{2\pi} \int_0^{2\pi} f(\theta) e^{-in\theta} d\theta \]
    For $1\leq p\leq 2$ and $\frac{1}{p}+\frac{1}{q}=1$, we have
    \[ ||Tf||_{L^q} \leq ||f||_{L^p} \]

    \end{lemma}
      \begin{proof}
      \uses{thm:riesz_interpolation} % Put any results used in the proof but not in the statement here
      % \leanok % uncomment if the lemma has been proven
      Observe that we may simply regard $T$ as an operator $L^1([0, 2\pi]) \to L^{\infty}(\bbz)$ since $L^2([0, 2\pi]) \subseteq L^1([0, 2\pi])$ (compact domain, bound with maximum), and $L^2(\bbz) \subseteq L^{\infty}(\bbz)$.\\
      Note that the claim corresponds (unless $q$ is infinity) to the inequality
      \[ \left(\sum_{n \in \bbz} |a_n|^q \right)^{1/q} \leq \left(\frac{1}{2\pi} \int_0^{2\pi} |f(\theta)|^p d\theta\right)^{1/p}\]
      For $p_0=2$ (thus $q_0=2$), this is Parseval's identity (see tsum\_sq\_fourierCoeff).\\
      For $p_1=1$ (thus $q_1=\infty$), we can check it directly. Since
      \[ a_n = \frac{1}{2\pi} \int_0^{2\pi} f(\theta) e^{-in\theta} d\theta \]
      we have
      \[ |a_n| \leq \frac{1}{2\pi} \int_0^{2\pi} |f(\theta) e^{-in\theta}| d\theta \leq \frac{1}{2\pi} \int_0^{2\pi} |f(\theta)| d\theta = ||f||_{L^1}  \]
      So
      \[ ||Tf||_{\infty} = \sup_n |a_n| \leq ||f||_{L^1}\]
      Applying Rietz-Thorin's theorem, we obtain that the claim holds whenever we can find a $t\in [0,1]$ such that
      \[\frac{1}{p} = \frac{1-t}{p_0} + \frac{t}{p_1} \qquad \frac{1}{q} = \frac{1-t}{q_0} + \frac{t}{q_1} \]
      Substituting $p_0=2$, $p_1=1$, $q_0=2$, $q_1=\infty$,
      \[\frac{1}{p} = \frac{1-t}{2} + t \qquad \frac{1}{q} = \frac{1-t}{2} \]
      Now
      \[ \frac{1}{p} = \frac{1+t}{2} \ \Rightarrow \ p = \frac{2}{1+t} \]
      which for $t\in [0,1]$ ranges from $1$ to $2$.\\
      Moreover, we have
      \[ \frac{1}{p} + \frac{1}{q} = \frac{1+t}{2} + \frac{1-t}{2} = 1 \]
      i.e., $p$ and $q$ are conjugate exponents.\\
      This completes the proof.
    \end{proof}

Now, we want to obtain a ``dual" inequality to the previous one. For this, we consider an operator $T':L^2(\bbz) \to L^2([0, 2\pi])$ in the opposite direction compared to the previous lemma
\[ T' (\{a_n\}_{n\in \bbz}) := \sum_{n=-\infty}^{\infty} a_n e^{in\theta} \]
The operator is defined on any $L^p(\bbz)$ for $p\leq 2$, since $L^p(\bbz) \subseteq L^2(\bbz)$.
Note that the target expression is indeed in $L^2([0, 2\pi])$ again.

\begin{lemma}
    \label{lem:hausdorff_young_dual}
    % \uses{}
    % \lean{}
    %\leanok
    For $1\leq p \leq 2$ and $q$ conjugate exponent to $p$, we have
    \[ ||T' \{a_n\}||_{L^q} \leq ||\{a_n\}||_{L^p} \]

    \end{lemma}
      \begin{proof}
      \uses{thm:riesz_interpolation} % Put any results used in the proof but not in the statement here
      % \leanok % uncomment if the lemma has been proven
      This is similar to the previous corollary. Parseval's identity gives the case $p_0=q_0=2$.\\
      For the case $p_1=1$, $q_1=\infty$, again
      \[ \left| \sum_{n\in \bbz} a_n e^{in\theta} \right| \leq \sum_{n\in \bbz} \left| a_n e^{in\theta} \right| = \sum_{n\in \bbz} |a_n| = ||\{a_n\}||_{L^1} \]
      i.e.
      \[ ||T'\{a_n\}||_{\infty} = \sup_{\theta \in [0, 2\pi]}  \left| \sum_{n\in \bbz} a_n e^{in\theta} \right| \leq ||\{a_n\}||_{L^1}  \]
      As before, applying Rietz-Thorin's interpolation theorem concludes the proof.
    \end{proof}

As a remark, if $f= T'\{a_n\}$, then the $\{a_n\}$ are the Fourier coefficients of $f$, yielding (when $p\neq 1$) the inequality
\[ \left( \frac{1}{2\pi} \int_0^{2\pi} |f(\theta)|^q d\theta \right)^{1/q} \leq \left( \sum_{n\in \bbz} |a_n|^p \right)^{1/p} \]

\subsection{Extending the Fourier transform}
The Rietz-Thorin interpolation theorem also allows us to extend the Fourier transform defined in the previous chapter to a bigger domain.\\
Let $V$ be a finite dimensional real inner product space and $E$ be a normed complex space.\\
As in Definition \ref{def:fourier-transform} define the Fourier transform on simple functions $f$ via the expression
\[ \mathcal{F} (f)(w) = \int e^{-2\pi i \langle v, w \rangle } f(v) \]
We have shown $\mathcal{F}$ extends to a bounded linear operator $L^1 \to L^{\infty}$, and to a bounded linear operator $L^2 \to L^2$ (see Theorem \ref{thm:fourier-is-l2-linear}).\\
By Rietz-Thorin interpolation theorem, it can be uniquely extended to bounded linear operators $L^p \to L^q$
whenever $1\leq p \leq 2$ and $q$ is conjugate to $p$.

\begin{comment} This subsection is to be potentially added later
\subsection{Young's Inequality for Convolutions}
Here, we work with $X=Y=\R^d$ with standard scalar product and the usual Lebesgue measure.\\

\begin{lemma}
    \label{lem:young_convolution}
    % \uses{}
    % \lean{}
    %\leanok
    For any $f\in L^p$ and $g \in L^r$, their convolution
    \[ (f * g) (x) = \int_{\R^d} f(x-y) g(y) dy \]
    is well-defined (i.e. the right-hand side is integrable for almost every $x$) and, if $\frac{1}{q} = \frac{1}{p} + \frac{1}{r} - 1$, we have
    \[ || f * g ||_{L^q} \leq ||f||_{L^p} ||g||_{L^r}\]

    \end{lemma}
      \begin{proof}
      \uses{thm:riesz_interpolation} % Put any results used in the proof but not in the statement here
      % \leanok % uncomment if the lemma has been proven
      \begin{itemize}
        \item{\textbf{Case 1:} assume $f$ and $g$ are simple functions. Fix $g$ and consider the operator $Tf = f * g$, which is linear by linearity of integrals.
        Let $M:= ||g||_{L^r}$. If $r'$ is the conjugate exponent to $r$, then by H\"older's inequality
        \[||T(f)||_{L^{\infty}} \]
        }
        TO BE COMPLETED
      \end{itemize}
    \end{proof}
\end{comment}

%%FILE blueprint/src/chapter/main.tex
% This is the main point of entry to the blueprint.
% Add chapters of the blueprint here.
% This file is not meant to be built. Build src/web.tex or src/print.text instead.

\title{Blueprint for the collaborative Formalization Seminar}

\author{}

\maketitle

\tableofcontents

\input{chapter/analysis_mathlib.tex}
\input{chapter/plancherel.tex}
\input{chapter/interpolation.tex}
\input{chapter/distributions.tex}


%%FILE blueprint/src/chapter/plancherel.tex
\chapter{Plancherel's Theorem}
\label{chap:plancherel}
\section{Basic Properties of the Fourier Transform}
In this section we record, mostly without proofs, basic statements about the Fourier transform on
$L^1$ functions. Most of these are already formalized in mathlib.

Let $(V,\mu),(W,\nu)$ be vector spaces over $\R$ with a $\sigma$-finite measure,
$E,F$ be normed spaces over $\C$ and let $L:V\times W\to \R$, $M:E\times F\to\C$
be bilinear maps.
\begin{definition}
  \label{def:fourier-transform}
  % \uses{}
  \lean{VectorFourier.fourierIntegral, Fourier.fourierIntegral, Real.fourierIntegral,
        Real.fourierIntegralInv}
  \leanok
  Let $f\in L^1(V,E)$. Its \emph{Fourier transform} (w.r.t. $L$) is the function
  $\mathcal Ff=\widehat f:W\to E$ given by
  $$\mathcal Ff(w):=\widehat f(w):=\int_V e^{-2\pi i L(v,w)}f(v)\,d\mu.$$
  The \emph{inverse Fourier transform} (w.r.t. $L$) is similarly defined as
  $$\mathcal F^{-1}f(w)=\check f(w):=\int_Ve^{2\pi iL(v,w)}f(v)\,d\mu.$$
\end{definition}

\begin{lemma}
  \label{lem:fourier-bounded}
  \uses{def:fourier-transform}
  \lean{VectorFourier.fourierIntegral_convergent_iff,
       VectorFourier.norm_fourierIntegral_le_integral_norm}
  \leanok
  Let $f\in L^1(V,E)$. Then its Fourier transform $\widehat f$ is well-defined and bounded.
  In particular, the Fourier transform defines a map $\mathcal F:L^1(V,E)\to L^\infty(V,E)$.
\end{lemma}
\begin{proof}
  \leanok
  Omitted.
\end{proof}

From now on assume that $V$ and $W$ are equipped with second-countable topologies such that
$L$ is continuous.
\begin{lemma}
  \label{lem:fourier-cont}
  \uses{def:fourier-transform}
  \lean{VectorFourier.fourierIntegral_continuous}
  \leanok
  Let $f\in L^1(V,E)$. Then $\widehat f$ is continuous.
\end{lemma}
\begin{proof}
  \leanok
  Omitted.
\end{proof}

\begin{lemma}[Multiplication formula]
\label{lem:fourier-multiplication}
\uses{lem:fourier-cont}
\lean{VectorFourier.integral_bilin_fourierIntegral_eq_flip}
\leanok
Let $f,g\in L^1(V,E)$. Then
$$\int_WM(\widehat f(w),g(w))\,d\nu(w)=\int_VM(f(v),\widehat g(v))\,d\mu(v).$$
\end{lemma}
\begin{proof}
  \leanok
  Omitted.
\end{proof}

\begin{lemma}
  \label{lem:fourier-prop}
  \uses{def:fourier-transform}
  \lean{VectorFourier.fourierIntegral_add,
        Fourier.fourierIntegral_comp_add_right,
        Fourier.fourierIntegral_smul_const,
        MeasureTheory.fourier_conj,
        MeasureTheory.fourier_convolution}
  \leanok
Lef $f,g\in L^1(V,E)$, $t\in\R$ and $a,b\in\C$. The Fourier transform satisfies the following elementary properties:
\begin{enumerate}
  \item[(i)] $\mathcal F(af+bg)=a\mathcal Ff+b\mathcal Fg$\hfill(Linearity)
  \item[(ii)] $\mathcal F(f(x-t)) = e^{-2\pi i ty}\mathcal F f(y)$\hfill (Shifting)
  \item[(iii)] $\mathcal F(f(tx)) = \frac1{|t|}\mathcal F f(\frac yt)$\hfill (Scaling)
  \item[(iv)] If $E$ admits a conjugation, then $\mathcal F(\overline{f(x)}) = \overline{\mathcal Ff(-y)}$\hfill (Conjugation)
  \item[(v)] Define the convolution of $f$ and $g$ w.r.t. a bilinear map $M:E\times E\to F$ as
             $$(f\ast_Mg)(w):=\int_VM(f(v),g(w-v))\,d\mu(v).$$ Then
             $\mathcal F(f\ast_M g) =M(\mathcal Ff,\mathcal Fg)$ \hfill (Convolution)
\end{enumerate}
\end{lemma}
\begin{proof}
  \leanok
  Omitted.
\end{proof}
From now on, let $V$ be a finite-dimensional inner product space.
We denote this product as ordinary multiplication, and the induced norm as $|\cdot|.$

We now study a family of functions which is useful for later proofs.
\begin{lemma}
\label{lem:fourier-gaussian}
\uses{def:fourier-transform}
\lean{fourierIntegral_gaussian_innerProductSpace'}
\leanok
Let $x\in V$ and $\delta>0$. Define the \emph{modulated Gaussian}
$$u_{x,\delta}(y):V\to\C,\quad y\mapsto e^{-\delta\pi|y|^2}e^{2\pi i x y}.$$
Its Fourier transform (w.r.t. the inner product) is given by
$$\widehat{u_{x,\delta}}(z)=\delta^{-n/2}e^{-\pi|x-y|^2/\delta}=:K_\delta(x-z).$$
\end{lemma}
\begin{proof}
  \uses{lem:fourier-prop}
  \leanok
  By choosing an orthonormal basis, wlog we may assume $V=\R^n$.
  First note $\widehat{u_{x,\delta}}(z-x)=\widehat{u_{0,\delta}}(z)$, so it is enough to consider $x=0$.
  Next, $$\widehat{u_{0,\delta}}(z)=\int_{\R^n}e^{-\pi\delta|y|^2-2\pi iyz}\,dy=
  \prod_{i=1}^n\int_{\R}e^{-\pi\delta y_i^2-2\pi iy_iz_i}dy_i\quad\text{and}\quad
  K_\delta(-z)=\prod_{i=1}^n\delta^{-1/2}e^{-\pi y_i^2/\delta},$$ hence we may assume $n=1$.
  The change of variables $w=\delta^{1/2}y+iz/\delta^{1/2}$ results in
  $$\widehat{u_{0,\delta}}(z)=\int_{\R} e^{-\pi\delta y^2-2\pi i y z}\,dy=
  \delta^{-1/2}e^{-\pi z^2/\delta}\int_{Im(w)=z/\delta^{1/2}}e^{-\pi w^2}\,dw.$$
  Contour integration along the rectangle with vertices $(\pm R,0),(\pm R,iz/\delta^{1/2})$, together
  with the bound $$\left|\int_{\pm R}^{\pm R+iz/\delta^{1/2}}e^{-\pi w^2}\,dw\right|\leq
  \frac{|z|}{\delta^{1/2}}\sup_{w\in[R,R+iz/\delta^{1/2}]}|e^{-\pi w^2}|=
  \frac{|z|}{\delta^{1/2}}e^{-\pi R^2}\xrightarrow{R\to\infty}0$$
  yields
  $$\int_{Im(w)=-z/\delta^{1/2}}e^{-\pi w^2}\,dw=
  \int_{\R} e^{-\pi w^2}\,dw=1,$$ finishing the proof.
\end{proof}

\begin{lemma}
\label{lem:weierstrass-kernel}
% \lean{}
\leanok
Let $K_\delta(v)=\delta^{-n/2}e^{-\pi|v|^2/\delta}$
as in \Cref{lem:fourier-gaussian}. This is a
\emph{good kernel}, called the \emph{Weierstrass kernel}, satisfying
$$\int_{V}K_\delta(x)\,dx= 1\quad\text{and}\quad
\int_{|x|>\eta}K_\delta(x)\,dx\xrightarrow{\delta\to0}0\quad\text{for all }\eta>0.$$
Furthermore, it satisfies the stronger bounds
$$K_\delta(x)\leq \delta^{-n/2}\quad\text{and}\quad K_\delta(x)\leq B\delta^{1/2}|x|^{-n-1}$$ for
some constant $B$ independent of $\delta$.
\end{lemma}
\begin{proof}
  \leanok
  By choosing an orthonormal basis, wlog we may assume $V=\R^n$. Then these are all straight-forward calculations:
  $$\int_{\R^n}e^{-\pi|x|^2/\delta}\,dx=\delta^{n/2}\int_{\R^n}e^{-\pi|x|^2}\,dx=\delta^{n/2}.$$
  $$\int_{|x|>\eta}\delta^{-n/2}e^{-\pi|x|^2/\delta}\,dx=\int_{|x|>\eta/\delta^{1/2}}e^{-\pi|x|^2}\xrightarrow{\delta\to0}0.$$
  The first upper bound is trivial. For the second one, consider for $r,z\geq0$ the inequality
  $$\Gamma(r+1)=\int_0^\infty e^{-y}y^r\,dy\geq\int_z^\infty e^{-y}y^r\,dy\geq z^r\int_z^\infty e^{-y}\,dy=z^re^{-z}.$$
  Applied to $z=\pi|x|^2/\delta$ and $r=(n+1)/2$, this gives
  $$|x|^{n+1}=\delta^{(n+1)/2}\frac{|x|^{n+1}}{\delta^{(n+1)/2}}\leq
  \underbrace{\frac{\Gamma((n+3)/2)}{\pi^{(n+1)/2}}}_{=:B}\delta^{(n+1)/2}e^{\pi|x|^2/\delta},$$
  which is equivalent to the second upper bound of the lemma.
\end{proof}

The following technical theorem is used in the proofs of both the inversion formula and Plancherel's theorem.
\begin{theorem}
\label{thm:kernel-approximation}
\uses{lem:weierstrass-kernel}
\lean{Real.tendsto_integral_gaussian_smul'}
\leanok
Let $f:V\to E$ be integrable. Let $K_\delta$ be the Weierstrass kernel from \Cref{lem:weierstrass-kernel},
or indeed any family of functions satisfying the conditions of \Cref{lem:weierstrass-kernel}. Then
$$(K_\delta\ast f)(x):=\int_V K_\delta(y)f(x-y)\,d\mu(y)\xrightarrow{\delta\to0}f(x)$$ in the $L^1$-norm.
If $f$ is continuous, the convergence also holds pointwise.
\end{theorem}
\begin{proof}
\leanok
Again we may assume $V=\R^n$. Consider the difference $$\Delta_\delta(x):=(K_\delta\ast f)(x)-f(x)=
\int_{\R^n}(f(x-y)-f(x))K_\delta(y)\,dy.$$ We prove $L^1$-convergence first:
Take $L^1$-norms and use Fubini's theorem to conclude
$$\|\Delta_\delta\|_1\leq\int_{\R^n}\|f(x-y)-f(x)\|_1K_\delta(y)\,dy.$$ For $\varepsilon>0$ find $\eta>0$
small enough so that $\|f(x-y)-f(x)\|_1<\varepsilon$ when $|y|\eta$. Thus
$$\|\Delta_\delta\|_1\leq\varepsilon+\int_{|y|>\eta}\|f(x-y)-f(x)\|K_\delta(y)\,dy\leq\varepsilon+2\|f\|_1\int_{\|y\|>\eta}K_\delta(y)\,dy.$$
By one of the properties in \Cref{lem:weierstrass-kernel}, we can choose $\delta$ small enough so that the second integral
is less than $\varepsilon$, which finishes the proof in this case.

Now assume that $f$ is continuous.
Let $d=\delta^{1/2}$ and shorten $g_\delta(x,y)=|f(x-y)-f(x)|K_\delta(y)$.Then
 $$|\Delta_\delta(x)|\leq\int_{0<|y|<d}g_\delta(x,y)\,dy+\sum_{k\in\mathbb N}\int_{2^kd<|y|<2^{k+1}d}g_\delta(x,y)\,dy.$$
 To bound these integrals, consider
 $$\varphi(r)=\frac1{r^n}\int_{|y|<r}|f(x-y)-f(x)|\,dy.$$
 It is easy to see that $\varphi$ is continuous, bounded, and approaches $0$ for $r\to0$, by continuity of $f$.
 Now $$\int_{0<|y|<d}g_\delta(x,y)\,dy\overset{(\ast)}\leq d^{-n}\int_{0<|y|<d}|f(x-y)-f(x)|\,dy=\varphi(d)$$ and
 $$\int_{2^kd<|y|<2^{k+1}d}g\,dy\overset{(\ast)}\leq
 \frac{Bd}{(2^kd)^{n+1}}\int_{2^kd<|y|<2^{k+1}d}|f(x-y)-f(x)|\,dy\leq2^{n-k}B\varphi(2^{k+1}d),$$
 where for the inequalities labeled $(\ast)$ we used the upper bounds from \Cref{lem:weierstrass-kernel}.
 Together, we find
 $$|\Delta_\delta(x)|\leq\varphi(d)+C\sum_{k\in\mathbb N}2^{-k}\varphi(2^{k+1}d)$$ for
 $C=\frac{2^n\Gamma((n+3)/2)}{\pi^{(n+1)/2}}$. Say $\varphi$ is bounded by $M\in\R$ and let $\varepsilon>0$.
 Take $N$ large enough such that $\sum_{k\geq N}2^{-k}<\varepsilon$. Choose $\delta$ small enough that
 $A(2^kd)<\varepsilon/N$ for all $k<N$. Then
 $$|\Delta_\delta(x)|\leq\varepsilon/N+(N-1)C\varepsilon/N+C\varepsilon M\leq \varepsilon C(M+1).$$
\end{proof}
\begin{remark}
  One can drop the continuity assumption and still get pointwise convergence almost everywhere. The proof stays the same,
  but one focuses on Lebesgue points of $f$. It takes slightly more work to argue that $\varphi$ behaves nicely,
  but the rest of the proof stays the same.
\end{remark}

\begin{theorem}[Inversion formula]
  \label{thm:fourier-inversion}
  \lean{MeasureTheory.Integrable.fourier_inversion, Continuous.fourier_inversion}
  Let $f:V\to E$ be integrable and continuous. Assume $\widehat f$ is integrable as well. Then
  $$\mathcal F^{-1}\mathcal F f=f.$$
  \leanok
\end{theorem}
\begin{proof}
  \uses{thm:kernel-approximation, lem:fourier-multiplication}
  \leanok
  Apply the multiplication formula \Cref{lem:fourier-multiplication} to $u_{x,\delta}$ and $f$, and conclude with
  \Cref{thm:kernel-approximation}.
\end{proof}

\begin{remark}
  Note that both assumptions are necessary, since $\mathcal F^{-1}\mathcal Ff$ is continuous, and
  only defined if $\mathcal Ff$ is integrable.
\end{remark}

\begin{theorem}[Inversion formula, $L^1$-version]
  \label{thm:fourier-inversion-L1}
  \leanok
  Let $f\in L^1(V,E)$. If $\widehat f\in L^1(V,E)$, then $\mathcal F^{-1}\mathcal Ff=f$.
\end{theorem}
\begin{proof}
  \leanok
  \uses{thm:kernel-approximation, lem:fourier-multiplication}
  Similar to \Cref{thm:fourier-inversion}.
\end{proof}

\section{\texorpdfstring{Plancherel's Theorem and the Fourier Transform on $L^2$}
                        {Plancherel's Theorem and the Fourier Transform on L2}}
Let $(V,\cdot)$ be a finite-dimensional inner product space over $\R$ and let $(E,\langle\cdot,\cdot\rangle)$ be an inner product space over $\C$.

\begin{theorem}[Plancherel's Theorem]
  \label{thm:plancherel}
  \lean{MeasureTheory.memℒp_fourierIntegral, MeasureTheory.snorm_fourierIntegral}
  \leanok
   Suppose that $f : V \to E$ is in $L^1(V,E)\cap L^2(V,E)$ and let $\widehat{f}$ be the Fourier transform of $f$. Then $\widehat{f},\check{f}\in L^2(V,E)$ and
  \[\|\widehat{f}\|_{L^2} = \|f\|_{L^2}=\|\check f\|_{L^2}.\]
   Suppose that $f : V \to E$ is in $L^1(V,E)\cap L^2(V,E)$ and let $\widehat{f}$ be the Fourier transform of $f$. Then $\widehat{f},\check{f}\in L^2(V,E)$ and
  \[\|\widehat{f}\|_{L^2} = \|f\|_{L^2}=\|\check f\|_{L^2}.\]
  \end{theorem}
    \begin{proof}
    \uses{thm:kernel-approximation, lem:fourier-multiplication, lem:fourier-gaussian, lem:fourier-prop}
    Let $g(x)=f(-x)$ and apply the multiplication formula \Cref{lem:fourier-multiplication}
    to $f\ast g$ and $u_{0,\delta}$: $$\int_V\widehat{f\ast g}\cdot u_{0,\delta}(x)\,dx=\int_V(f\ast g)(x)K_\delta(-x)\,dx
    \overset{\delta\to0}\to(f\ast g)(0)=\int_V \langle f(x),f(x)\rangle\,dx=\| f\|_2^2$$ by \Cref{thm:kernel-approximation}.
    On the other hand, by \Cref{lem:fourier-prop} the left hand side simplifies to
    $$\int_V|\widehat f(x)|^2e^{-\delta\pi|x|^2}\,dx\xrightarrow{\delta\to0}\|\widehat f\|_2^2$$ by dominated convergence.

    Since $\check f(x)=\widehat f(-x)$, the corresponding statements for $\check f$ follow immediately from the ones for $\widehat f$.
\end{proof}

We now want to extend the Fourier transform to $L^2(V,E)$. For this, take a sequence of functions $(f_n)_n\subset L^1(V,E)\cap L^2(V,E)$
such that $f_n\xrightarrow[L^2]{}f$. Such sequences exist:
\begin{lemma}
  \label{lem:L12-dense}
  %\uses{MeasureTheory.Memℒp.exists_hasCompactSupport_snorm_sub_le, Mathlib.MeasureTheory.Function.SimpleFuncDenseLp}
  \lean{}
  %\leanok
  $L^1(V,E)\cap L^2(V,E)$ is dense in $L^2(V,E)$.
\end{lemma}
\begin{proof}
  %\leanok
  It is well-known that the space of compactly supported continuous functions is dense in every $L^p(V,E)$.
  Since those are contained in $L^1(V,E)\cap L^2(V,E)$, the claim follows.
\end{proof}

Let $f\in L^2(V,E)$. Plancherel's theorem lets us now approximate a potential $\widehat f$:
\begin{lemma}
  \label{lem:fourier12-cauchy}
  \uses{lem:L12-dense,}
  \lean{}
  % \leanok
  Let $f\in L^2(V,E)$ and $(f_n)_n\subset L^1(V,E)\cap L^2(V,E)$ a sequence with $f_n\xrightarrow[L^2]{}f$. Then $(\widehat f_n)_n$ is a Cauchy sequence,
  hence converges in $L^2(V,E)$.
\end{lemma}
\begin{proof}
  % \leanok
  \uses{thm:plancherel, lem:fourier-prop}
  $$\|\widehat f_n-\widehat f_m\|_2=\|\widehat{f_n-f_m}\|_2\overset{\text{Plancherel}}=\|f_n-f_m\|_2$$ goes to $0$ for $n,m$ large, as $(f_n)_n$ is
  convergent, hence Cauchy. Since $L^2(V,E)$ is complete, $(\widehat f_n)_n$ converges.
\end{proof}
\begin{definition}
  \label{def:fourier-L2}
  \uses{lem:fourier12-cauchy}
  \lean{}
  % \leanok
  Let $f\in L^2(V,E)$ and take a sequence $(f_n)_n\subset L^1(V,E)\cap L^2(V,E)$ with $f_n\xrightarrow[L^2]{}f$. Set
  $$\mathcal Ff:=\widehat f:=\lim_{n\to\infty}\widehat{f_n},$$ the limit taken in the $L^2$-sense.
\end{definition}
\begin{lemma}
  \label{lem:fourier2-welldef}
  \uses{def:fourier-L2, lem:fourier12-cauchy}
  \lean{}
  % \leanok
  This is well-defined: By \Cref{lem:fourier12-cauchy}, the limit exists. Further it does not depend
  on the choice of sequence $(f_n)_n$. If $f\in L^1(V,E)\cap L^2(V,E)$, this definition agrees with
  the Fourier transform on $L^1(V,E)$.
\end{lemma}
\begin{proof}
  % \leanok
  \uses{thm:plancherel}
  Let $(g_n)_n$ be another sequence approximating $f$. Then
  $$\|\widehat f_n-\widehat g_n\|_2=\|f_n-g_n\|\leq\|f_n-f\|+\|g_n-g\|\to0.$$
  If $f\in L^1(V,E)\cap L^2(V,E)$, we can choose the constant sequence $(f_n)_n=(f)_n$.
\end{proof}

\begin{definition}
  \label{def:invFourier-L2}
  \uses{lem:fourier12-cauchy, thm:plancherel}
  \lean{}
  % \leanok
  Define analogously $\mathcal F^{-1}f:=\check f:=\lim_{n\to\infty}\check f_n$, if $f_n\xrightarrow[L^2]{}f\in L^2(V,E)$
  with $(f_n)_n\subset L^1(V,E)\cap L^2(V,E)$. By the same arguments as above, this is well-defined.
\end{definition}

\begin{corollary}
  \label{thm:fourier2-properties}
  \uses{def:fourier-L2, def:invFourier-L2, lem:fourier2-welldef}
  \lean{}
  % \leanok
  Plancherel's Theorem, the inversion formula, and the properties of \Cref{lem:fourier-prop} hold
  for the Fourier transform on $L^2(V,E)$ as well.
\end{corollary}
\begin{proof}
  \uses{thm:plancherel, thm:fourier-inversion, lem:fourier-prop}
  % \leanok
  All of these follow immediately from the definition and the observation, that all operations (norms, sums, conjugation, $\ldots$) are continuous.
  For example, let $f\in L^2(V,E)$ and take an approximating sequence $(f_n)_n$ as before. Then
  $$\|\widehat f\|_2=\|\lim_{n\to\infty}\widehat f_n\|_2=\|\lim_{n\to\infty}\|\widehat f_n\|_2=\lim_{n\to\infty}\|f_n\|_2
  =\|\lim_{n\to\infty}f_n\|_2=\|f\|_2.$$
\end{proof}

\begin{corollary}
  \label{thm:fourier-is-l2-linear}
  \uses{def:fourier-L2, lem:fourier2-welldef}
  \lean{MeasureTheory.fourierIntegralL2}
  \leanok
  The Fourier transform induces a continuous linear map $L^2(V,E) \to L^2(V,E)$.
  \end{corollary}
  \begin{proof}
    \uses{thm:fourier2-properties}
    % \leanok % uncomment if the lemma has been proven
    This follows immediately from \Cref{thm:fourier2-properties}: Linearity from the $L^2$-version of
    \Cref{lem:fourier-prop}, and continuity and well-definedness from the $L^2$-version of Plancherel's theorem.
\end{proof}

